\chapter{Epicureans}

\lettrine{A}{s} Fouquet was giving, or appearing to give, all his attention to the brilliant illuminations, the languishing music of the violins and hautboys, the sparkling sheaves of the artificial fires, which, inflaming the heavens with glowing reflections, marked behind the trees the dark profile of the donjon of Vincennes; as, we say, the superintendent was smiling on the ladies and the poets, the fête was every whit as gay as usual; and Vatel, whose restless, even jealous look, earnestly consulted the aspect of Fouquet, did not appear dissatisfied with the welcome given to the ordering of the evening's entertainment. The fireworks over, the company dispersed about the gardens and beneath the marble porticoes with the delightful liberty which reveals in the master of the house so much forgetfulness of greatness, so much courteous hospitality, so much magnificent carelessness. The poets wandered about, arm in arm, through the groves; some reclined upon beds of moss, to the great damage of velvet clothes and curled heads, into which little dried leaves and blades of grass insinuated themselves. The ladies, in small numbers, listened to the songs of the singers and the verses of the poets; others listened to the prose, spoken with much art, by men who were neither actors nor poets, but to whom youth and solitude gave an unaccustomed eloquence, which appeared to them better than everything else in the world. <Why,> said La Fontaine, <does not our master Epicurus descend into the garden? Epicurus never abandoned his pupils; the master is wrong.>

<Monsieur,> said Conrart, <you yourself are in the wrong persisting in decorating yourself with the name of an Epicurean; indeed, nothing here reminds me of the doctrine of the philosopher of Gargetta.>

<Bah!> said La Fontaine, <is it not written that Epicurus purchased a large garden and lived in it tranquilly with his friends?>

<That is true.>

<Well, has not M.~Fouquet purchased a large garden at Saint-Mande, and do we not live here very tranquilly with him and his friends?>

<Yes, without doubt; unfortunately it is neither the garden nor the friends which constitute the resemblance. Now, what likeness is there between the doctrine of Epicurus and that of M.~Fouquet?>

<This—pleasure gives happiness.>

<Next?>

<Well, I do not think we ought to consider ourselves unfortunate, for my part, at least. A good repast—vin de Joigny, which they have the delicacy to go and fetch for me from my favourite cabaret—not one impertinence heard during a supper an hour long, in spite of the presence of ten millionaires and twenty poets.>

<I stop you there. You mentioned vin de Joigny, and a good repast; do you persist in that?>

<I persist,—anteco, as they say at Port Royal.>

<Then please to recollect that the great Epicurus lived, and made his pupils live, upon bread, vegetables, and water.>

<That is not certain,> said La Fontaine; <and you appear to me to be confounding Epicurus with Pythagoras, my dear Conrart.>

<Remember, likewise, that the ancient philosopher was rather a bad friend of the gods and the magistrates.>

<Oh! that is what I will not admit,> replied La Fontaine. <Epicurus was like M.~Fouquet.>

<Do not compare him to monsieur le surintendant,> said Conrart, in an agitated voice, <or you would accredit the reports which are circulating concerning him and us.>

<What reports?>

<That we are bad Frenchmen, lukewarm with regard to the king, deaf to the law.>

<I return, then, to my text,> said La Fontaine. <Listen, Conrart, this is the morality of Epicurus, whom, besides, I consider, if I must tell you so, as a myth. Antiquity is mostly mythical. Jupiter, if we give a little attention to it, is life. Alcides is strength. The words are there to bear me out; Zeus, that is, zen, to live. Alcides, that is, alce, vigour. Well, Epicurus, that is mild watchfulness, that is protection; now who watches better over the state, or who protects individuals better than M.~Fouquet does?>

<You talk etymology and not morality; I say that we modern Epicureans are indifferent citizens.>

<Oh!> cried La Fontaine, <if we become bad citizens, it is not through following the maxims of our master. Listen to one of his principal aphorisms.>

<I—will.>

<Pray for good leaders.>

<Well?>

<Well! what does M.~Fouquet say to us every day? <When shall we be governed?> Does he say so? Come, Conrart, be frank.>

<He says so, that is true.>

<Well, that is a doctrine of Epicurus.>

<Yes; but that is a little seditious, observe.>

<What! seditious to wish to be governed by good heads or leaders?>

<Certainly, when those who govern are bad.>

<Patience, I have a reply for all.>

<Even for what I have just said to you?>

<Listen! would you submit to those who govern ill? Oh! it is written: \foreign{Cacos politeuousi.} You grant me the text?>

<\swear{Pardieu!} I think so. Do you know, you speak Greek as well as Aesop did, my dear La Fontaine.>

<Is there any wickedness in that, my dear Conrart?>

<God forbid I should say so.>

<Then let us return to M.~Fouquet. What did he repeat to us all the day? Was it not this? <What a cuistre is that Mazarin! what an ass! what a leech! We must, however, submit to that fellow.> Now, Conrart, did he say so, or did he not?>

<I confess that he said it, and even perhaps too often.>

<Like Epicurus, my friend, still like Epicurus; I repeat, we are Epicureans, and that is very amusing.>

<Yes; but I am afraid there will rise up, by the side of us, a sect like that of Epictetus; you know him well; the philosopher of Hierapolis, he who called bread luxury, vegetables prodigality, and clear water drunkenness; he who, being beaten by his master, said to him, grumbling a little it is true, but without being angry, <I will lay a wager you have broken my leg!>—and who won his wager.>

<He was a goose, that fellow Epictetus.>

<Granted, but he might easily become the fashion by only changing his name into that of Colbert.>

<Bah!> replied La Fontaine, <that is impossible. Never will you find Colbert in Epictetus.>

<You are right, I shall find—Coluber there, at the most.>

<Ah! you are beaten, Conrart; you are reduced to a play upon words. M.~ Arnaud pretends that I have no logic; I have more than M.~Nicole.>

<Yes,> replied Conrart, <you have logic, but you are a Jansenist.>

This peroration was hailed with a boisterous shout of laughter; by degrees the promenaders had been attracted by the exclamations of the two disputants around the arbour under which they were arguing. The discussion had been religiously listened to, and Fouquet himself, scarcely able to suppress his laughter, had given an example of moderation. But with the denouement of the scene he threw off all restraint, and laughed aloud. Everybody laughed as he did, and the two philosophers were saluted with unanimous felicitations. La Fontaine, however, was declared conqueror, on account of his profound erudition and his irrefragable logic. Conrart obtained the compensation due to an unsuccessful combatant; he was praised for the loyalty of his intentions, and the purity of his conscience.

At the moment when this jollity was manifesting itself by the most lively demonstrations, when the ladies were reproaching the two adversaries with not having admitted women into the system of Epicurean happiness, Gourville was seen hastening from the other end of the garden, approaching Fouquet, and detaching him, by his presence alone, from the group. The superintendent preserved on his face the smile and character of carelessness; but scarcely was he out of sight than he threw off the mask.

<Well!> said he, eagerly, <where is Pélisson! What is he doing?>

<Pélisson has returned from Paris.>

<Has he brought back the prisoners?>

<He has not even seen the concierge of the prison.>

<What! did he not tell him he came from me?>

<He told him so, but the concierge sent him this reply: <If any one came to me from M.~Fouquet, he would have a letter from M.~Fouquet.>>

<Oh!> cried the latter, <if a letter is all he wants\longdash>

<It is useless, monsieur!> said Pélisson, showing himself at the corner of the little wood, <useless! Go yourself, and speak in your own name.>

<You are right. I will go in, as if to work; let the horses remain harnessed, Pélisson. Entertain my friends, Gourville.>

<One last word of advice, monseigneur,> replied the latter.

<Speak, Gourville.>

<Do not go to the concierge save at the last minute; it is brave, but it is not wise. Excuse me, Monsieur Pélisson, if I am not of the same opinion as you; but take my advice, monseigneur, send again a message to this concierge,—he is a worthy man, but do not carry it yourself.>

<I will think of it,> said Fouquet; <besides, we have all the night before us.>

<Do not reckon too much on time; were the hours we have twice as many as they are, they would not be too much,> replied Pélisson; <it is never a fault to arrive too soon.>

<Adieu!> said the superintendent; <come with me, Pélisson. Gourville, I commend my guests to your care.> And he set off. The Epicureans did not perceive that the head of the school had left them; the violins continued playing all night long. 